%!TeX root = main.tex
%!encoding = UTF-8 Unicode

\subsection{Design Constraints and Measurement Requirements}
\label{sec:CurrentRequirements}

The following requirements guide the design of the dual-channel active current probe for commutation-current measurements on discrete SiC MOSFETs.
The measurement criteria discussed for the voltage probe in Section~\ref{sec:VoltageRequirements} also apply to the current probe.
The general error mechanisms described in Section~\ref{sec:fundamental_measurement_errors} must likewise be considered.
Current sensing introduces an additional constraint: the insertion resistance must be low enough to preserve the commutation-loop damping.
The discussion below therefore briefly addresses both measurement accuracy and the probe's influence on the switching transient.

\subsubsection*{Small Signal Response}
The current probe must support both conventional switching-current measurements and the analysis of the EMC-relevant,high-frequency range.
In the measurements considered here, the switching-current spectrum approaches the noise floor at approximately \SI{120}{MHz}.
The low-pass channel must therefore have a \SI{-3}{dB} bandwidth above this frequency, while the high-pass channel must retain a rising magnitude response up to \SI{500}{MHz}, the upper limit for EMC-relevant components considered here (Section~\ref{sec:VoltageRequirements}).

Conventional current sensors generally fall short of the high-pass target: typical usable bandwidths are approximately \SIrange{20}{50}{MHz} for commercial Rogowski coils, \SI{10}{MHz} to \SI{100}{MHz} for magnetoresistive sensors, and \SIrange{1}{10}{MHz} for Hall-effect sensors~\cite{xin_review_2022}.
Custom Rogowski coils with very few turns can reach higher frequencies, but their response is sensitive to position and orientation, requiring careful calibration and controlled placement.
A shunt-based approach is therefore selected to support high-bandwidth measurements with robust integration into the commutation cell.

\subsubsection*{Large Signal Response}
The current probe must capture switching-current peaks, which can substantially exceed the device's nominal current rating, particularly at high temperatures.
Unlike the rated blocking voltage used to specify the voltage-probe range, the nominal current does not provide an upper bound for the current-probe range.
The required range is therefore based on the largest observed peak.

For the investigated \SI{60}{A} SiC MOSFET in a D2PAK package, measurements under worst-case high-temperature conditions showed peak currents up to \SI{140}{A} (Section~\ref{sec:Results}).
Accordingly, the probe's maximum current rating was set to \SI{150}{A}, providing margin above the observed peak.

Because current and voltage transients are governed by the same commutation-loop LC resonance, the current probe must resolve rise times comparable to the lower-bound estimate for the voltage probe (Section~\ref{sec:VoltageRequirements}).

\subsubsection*{Loading and Measurement Intrusion}
The shunt must not materially change the behavior of the switching transient.
The dominant resonance involves the commutation-loop inductance and the DUT's parasitic capacitances.
Its quality factor is an EMC-relevant characteristic, so additional damping from the measurement would alter the behavior under investigation.
As a design rule, the inserted resistance is limited to one tenth of the intrinsic damping resistance of the commutation cell.

Although recent research points to additional damping mechanisms \cite{GaNRoss}, conventional analyses often treat the switching device's on-state resistance as the dominant contribution to commutation-loop damping.
Values range from \SIrange{20}{50}{m\ohm} for the considered \SI{750}{V} and \SI{1200}{V} SiC MOSFETs in D2PAK packages~\cite{robert_bosch_gmbh_datasheet_2023,robert_bosch_gmbh_datasheet_2023-1}.
The sense resistance must therefore not exceed \SI{2}{m\ohm}.

Low insertion resistance also limits the bandwidth of an uncompensated shunt, motivating active compensation.

\subsubsection*{Parasitic Coupling}
The coupling mechanisms discussed in Section~\ref{sec:VoltageRequirements} also apply to the current probe.
Coupling from the switching node is particularly critical because the sensing circuit is placed nearby and operates with higher gain than the voltage probe.
Layout and shielding must therefore minimize coupling into the sensitive amplification path.

\subsubsection*{Transmission Line Effects}
The short-line criterion derived in Section~\ref{sec:VoltageRequirements} also applies to the current-sensing circuit.
At \SI{500}{MHz}, the same FR4 assumption limits non-\SI{50}{\ohm} connections to \SI{2.9}{cm}.
This constraint applies to the interconnections between the shunt element, compensation amplifier, and filter stage.
These stages must therefore be integrated closely rather than connected through unmatched cables.

\subsubsection*{Output Interface and Signal Routing}
As in the voltage-probe concept, a \SI{50}{\ohm} output interface is required for the low-pass and high-pass signal paths.
The milliohm-range shunt cannot directly drive this interface without additional signal conditioning.
An active stage must therefore provide both the required output level and the defined interface impedance.

\subsubsection*{Isolation for Combined Measurements}
For standalone current measurements, the probe does not require galvanic isolation.
Isolation becomes necessary only when commutation current and LISN disturbance voltage are acquired simultaneously: the oscilloscope ground is then fixed to the aluminium base plate and cannot serve as the shunt reference.
An optically isolated voltage probe can then provide the required separation while retaining the \SI{50}{\ohm} signal interface.
\\
\\
The design requirements are summarized in \tref{tab:CurrentProbeRequirements}.

\begin{table}[htbp]
\centering
\caption{Summary of the dual-channel active current probe requirements.}
\label{tab:CurrentProbeRequirements}
\renewcommand{\arraystretch}{1.25}
\begin{tabularx}{0.8\textwidth}{>{\raggedright\arraybackslash}p{0.28\textwidth} >{\centering\arraybackslash}X}
    \toprule
    \textbf{Requirement} & \textbf{Design limit} \\
    \midrule
    Low-pass response & $\mi{f}{\SI{3}{dB}}>\SI{120}{MHz}$ \\
    High-pass response & $\mi{f}{interest,max} = \SI{500}{MHz}$ \\
    Peak current & $I_{\mathrm{max}}= \SI{150}{A}$ \\
    Sense resistance & $R_{\mathrm{shunt}}\leq\SI{2}{m\Omega}$ \\
    Internal connections & $l_{\max}\leq\SI{2.9}{cm}$ \\
    Output interface & \SI{50}{\ohm} \\
    Galvanic isolation & Required for combined current and LISN measurements \\
    \bottomrule
\end{tabularx}
\end{table}

\FloatBarrier